\section{Adic rings}
\label{section:0.7}

\subsection{Admissible rings}
\label{subsection:0.7.1}

\begin{env}[7.1.1]
\label{0.7.1.1}
Recall that in a topological ring $A$ (not necessarily separated), an
element $x$ is said to be topologically nilpotent if $0$ is a limit of
the sequence $(x^n)_{n\geq0}$. A topological ring $A$ is said to be
linearly topologized if there exists a fundamental system of neighbourhoods of
$0$ in $A$ consisting of ideals (necessarily open).
\end{env}

\begin{definition}[7.1.2]
\label{0.7.1.2}
In a linearly topologized ring $A$, an ideal $\mathfrak J$ is said to be
an ideal of definition if $\mathfrak J$ is open and if, for every
neighbourhood $V$ of $0$, there exists an integer $n>0$ such

\oldpage[0\textsubscript{I}]{61}

that $\mathfrak J^n\subset V$ (which is expressed, by abuse of language, by saying
that the sequence $(\mathfrak J^n)$ tends to $0$). A linearly
topologized ring $A$ is said to be preadmissible if there exists in $A$ an ideal of definition;
$A$ is said to be admissible if it is preadmissible and if moreover it is
separated and complete.
\end{definition}

It is clear that if $\mathfrak J$ is an ideal of definition and $\mathfrak L$ an
open ideal of $A$, then $\mathfrak J\cap\mathfrak L$ is again an ideal of
definition; the ideals of definition of a preadmissible ring $A$ therefore form
a fundamental system of neighbourhoods of $0$.

\begin{lemma}[7.1.3]
\label{0.7.1.3}
Let $A$ be a linearly topologized ring.
\begin{enumerate}
  \item[{\upshape(i)}] For $x\in A$ to be topologically nilpotent, it is
    necessary and sufficient that for every open ideal $\mathfrak J$ of $A$, the canonical
    image of $x$ in $A/\mathfrak J$ be nilpotent. The set
    $\mathfrak T$ of topologically nilpotent elements of $A$ is an ideal.
  \item[{\upshape(ii)}] Suppose moreover that $A$ is preadmissible, and let
    $\mathfrak J$ be an ideal of definition of $A$. For $x\in A$ to be
    topologically nilpotent, it is necessary and sufficient that its canonical image in
    $A/\mathfrak J$ be nilpotent; the ideal $\mathfrak T$ is the inverse
    image of the nilradical of $A/\mathfrak J$ and is therefore open.
\end{enumerate}
\end{lemma}

(i) follows immediately from the definitions. To prove (ii), it suffices to
note that for every neighbourhood $V$ of $0$ in $A$, there exists $n>0$ such that
$\mathfrak J^n\subset V$; if $x\in A$ is such that $x^m\in\mathfrak J$, we have
$x^{mq}\in V$ for $q\geq n$, so $x$ is topologically nilpotent.

\begin{proposition}[7.1.4]
\label{0.7.1.4}
Let $A$ be a preadmissible ring and $\mathfrak J$ an ideal of definition of
$A$.
\begin{enumerate}
  \item[{\upshape(i)}] For an ideal $\mathfrak J'$ of $A$ to be contained
    in an ideal of definition, it is necessary and sufficient that there exist an integer
    $n>0$ such that ${\mathfrak J'}^n\subset\mathfrak J$.
  \item[{\upshape(ii)}] For an $x\in A$ to be contained in an ideal of
    definition, it is necessary and sufficient that it be topologically nilpotent.
\end{enumerate}
\end{proposition}

(i) If ${\mathfrak J'}^n\subset\mathfrak J$, then for every open neighbourhood $V$ of
$0$ in $A$, there exists $m$ such that $\mathfrak J^m\subset V$, hence
${\mathfrak J'}^{mn}\subset V$.

(ii) The condition is obviously necessary; it is sufficient, because if it
is satisfied, there exists $n$ such that $x^n\in\mathfrak J$, so
$\mathfrak J'=\mathfrak J+Ax$ is an ideal of definition, since it is open
and ${\mathfrak J'}^n\subset\mathfrak J$.

\begin{corollary}[7.1.5]
\label{0.7.1.5}
In a preadmissible ring $A$, an open prime ideal contains all the
ideals of definition.
\end{corollary}

\begin{corollary}[7.1.6]
\label{0.7.1.6}
The notation and hypotheses being those of (7.1.4), the following properties
of an ideal $\mathfrak J_0$ of $A$ are equivalent:
\begin{enumerate}
  \item[$a)$] $\mathfrak J_0$ is the largest ideal of definition of $A$;
  \item[$b)$] $\mathfrak J_0$ is a maximal ideal of definition;
  \item[$c)$] $\mathfrak J_0$ is an ideal of definition such that the ring
    $A/\mathfrak J_0$ is reduced.
\end{enumerate}
For there to exist an ideal $\mathfrak J_0$ having these properties, it is necessary and
sufficient that the nilradical of $A/\mathfrak J$ be nilpotent; $\mathfrak J_0$
is then equal to the ideal $\mathfrak T$ of topologically nilpotent elements
of $A$.
\end{corollary}

It is clear that $a)$ implies $b)$, and $b)$ implies $c)$ by virtue of
(7.1.4, (ii)) and (7.1.3, (ii)); for the same reason, $c)$ implies $a)$. The
last assertion follows from (7.1.4, (i)) and (7.1.3, (ii)).

When $\mathfrak T/\mathfrak J$, the nilradical of $A/\mathfrak J$, is
nilpotent, we denote by $A_{\mathrm{red}}$ the (reduced) quotient ring
$A/\mathfrak T$.

\oldpage[0\textsubscript{I}]{62}

\begin{corollary}[7.1.7]
\label{0.7.1.7}
A Noetherian preadmissible ring has a largest ideal of definition.
\end{corollary}

\begin{corollary}[7.1.8]
\label{0.7.1.8}
If a preadmissible ring $A$ is such that, for one ideal of definition
$\mathfrak J$, the powers $\mathfrak J^n$ ($n>0$) form a fundamental
system of neighbourhoods of $0$, then the same holds for the powers
${\mathfrak J'}^n$ of every ideal of definition $\mathfrak J'$ of $A$.
\end{corollary}

\begin{definition}[7.1.9]
\label{0.7.1.9}
A preadmissible ring $A$ is said to be preadic if there exists an ideal of
definition $\mathfrak J$ of $A$ such that the $\mathfrak J^n$ form a
fundamental system of neighbourhoods of $0$ in $A$ (or, what amounts to the same, such that
the $\mathfrak J^n$ are open). An adic ring is a separated and complete
preadic ring.
\end{definition}

If $\mathfrak J$ is an ideal of definition of a preadic (resp.
adic) ring $A$, we also say that $A$ is a
\emph{$\mathfrak J$-preadic} (resp. \emph{$\mathfrak J$-adic}) ring, and that its
topology is the \emph{$\mathfrak J$-preadic} (resp.
\emph{$\mathfrak J$-adic}) topology. More generally, if $M$ is an $A$-module, the
topology on $M$ having as a fundamental system of neighbourhoods of $0$ the
submodules $\mathfrak J^nM$ is called the
\emph{$\mathfrak J$-preadic} (resp. \emph{$\mathfrak J$-adic}) topology. By virtue
of (7.1.8), these topologies are independent of the choice of the ideal of definition
$\mathfrak J$.

\begin{proposition}[7.1.10]
\label{0.7.1.10}
Let $A$ be an admissible ring and $\mathfrak J$ an ideal of definition of $A$.
Then $\mathfrak J$ is contained in the radical of $A$.
\end{proposition}

This statement is equivalent to any one of the following corollaries:

\begin{corollary}[7.1.11]
\label{0.7.1.11}
For every $x\in\mathfrak J$, $1+x$ is invertible in $A$.
\end{corollary}

\begin{corollary}[7.1.12]
\label{0.7.1.12}
For $f\in A$ to be invertible in $A$, it is necessary and sufficient that its canonical
image in $A/\mathfrak J$ be invertible in $A/\mathfrak J$.
\end{corollary}

\begin{corollary}[7.1.13]
\label{0.7.1.13}
For every $A$-module $M$ of finite type, the relation $M=\mathfrak JM$
(equivalent to $M\otimes_A(A/\mathfrak J)=0$) implies $M=0$.
\end{corollary}

\begin{corollary}[7.1.14]
\label{0.7.1.14}
Let $u:M\to N$ be a homomorphism of $A$-modules, $N$ being of finite type;
for $u$ to be surjective, it is necessary and sufficient that
$u\otimes1:M\otimes_A(A/\mathfrak J)\to N\otimes_A(A/\mathfrak J)$ be so.
\end{corollary}

Indeed, the equivalence of (7.1.10) and (7.1.11) follows from Bourbaki,
\emph{Alg.}, chap.~VIII, \textsection6, n\textsuperscript{o}~3, th.~1, and
that of (7.1.10) and (7.1.13) from \emph{loc. cit.}, th.~2; the fact that
(7.1.10) implies (7.1.14) follows from \emph{loc. cit.}, cor.~4 of prop.~6;
on the other hand, (7.1.14) implies (7.1.13) by applying it to the zero homomorphism.
Finally, (7.1.10) implies that if $f$ is invertible in $A/\mathfrak J$, $f$
is not contained in any maximal ideal of $A$, so $f$ is invertible in
$A$; in other words, (7.1.10) implies (7.1.12); and conversely, (7.1.12)
implies (7.1.11).

Everything therefore comes down to proving (7.1.11). Now, since $A$ is separated and complete and
the sequence $(\mathfrak J^n)$ tends to $0$, it is immediate that the series
$\sum_{n=0}^{\infty}(-1)^n x^n$ is convergent in $A$ and that, if $y$ is
its sum, we have $y(1+x)=1$.

\subsection{Adic rings and projective limits}
\label{subsection:0.7.2}

\begin{env}[7.2.1]
\label{0.7.2.1}
Every projective limit of discrete rings is obviously a linearly
topologized, separated and complete ring. Conversely, let $A$ be a linearly
topologized ring, and let $(\mathfrak J_\lambda)$ be a fundamental system
of open neighbourhoods of $0$ in $A$ consisting

\oldpage[0\textsubscript{I}]{63}

of ideals. The canonical maps
$\varphi_\lambda:A\to A/\mathfrak J_\lambda$ form a projective system of
continuous maps and therefore define a continuous map
$\varphi:A\to\varprojlim A/\mathfrak J_\lambda$; if $A$ is separated,
$\varphi$ is a topological isomorphism of $A$ onto an
everywhere dense subring of $\varprojlim A/\mathfrak J_\lambda$; if,
moreover, $A$ is complete, $\varphi$ is a topological isomorphism of $A$ onto
$\varprojlim A/\mathfrak J_\lambda$.
\end{env}

\begin{lemma}[7.2.2]
\label{0.7.2.2}
For a linearly topologized ring to be admissible, it is necessary and sufficient
that it be isomorphic to a projective limit
$A=\varprojlim A_\lambda$, where $(A_\lambda,u_{\lambda\mu})$ is a
projective system of discrete rings having as index set a filtered ordered
set $L$ (for $\leq$) which has a smallest element, denoted $0$, and satisfies
the following conditions: 1\textsuperscript{o} the
$u_\lambda:A\to A_\lambda$ are surjective; 2\textsuperscript{o} the kernel
$\mathfrak J_\lambda$ of $u_{0\lambda}:A_\lambda\to A_0$ is nilpotent.
When this is so, the kernel $\mathfrak J$ of $u_0:A\to A_0$ is equal to
$\varprojlim\mathfrak J_\lambda$.
\end{lemma}

The necessity of the condition follows from (7.2.1), by taking for
$(\mathfrak J_\lambda)$ a fundamental system of neighbourhoods of $0$ consisting
of ideals of definition contained in one of them, $\mathfrak J_0$, and applying
(7.1.4, (i)). The converse follows from the definition of a projective limit and
from (7.1.2), and the last assertion is immediate.

\begin{env}[7.2.3]
\label{0.7.2.3}
Let $A$ be an admissible topological ring and $\mathfrak J$ an ideal of $A$
contained in an ideal of definition (in other words (7.1.4) such that
$(\mathfrak J^n)$ tends to $0$); we can consider on $A$ the ring topology
having as a fundamental system of neighbourhoods of $0$ the powers
$\mathfrak J^n$ ($n>0$); we shall again call it the
\emph{$\mathfrak J$-preadic} topology. The hypothesis that $A$ is admissible implies that
$\bigcap_n\mathfrak J^n=0$, so the $\mathfrak J$-preadic topology on $A$
is separated; let
$\widehat A=\varprojlim A/\mathfrak J^n$ be the completion of $A$ for this topology
(where the $A/\mathfrak J^n$ are endowed with the discrete topology), and denote
by $u$ the ring homomorphism $A\to\widehat A$ (not necessarily
continuous) that is the projective limit of the sequence of homomorphisms
$u_n:A\to A/\mathfrak J^n$. On the other hand, the $\mathfrak J$-preadic topology
on $A$ is finer than the given topology $\mathcal T$ on $A$; since $A$
is separated and complete for $\mathcal T$, we can extend by continuity
the identity map from $A$ (endowed with the $\mathfrak J$-preadic topology)
to $A$ endowed with $\mathcal T$; this gives a continuous map
$v:\widehat A\to A$.
\end{env}

\begin{proposition}[7.2.4]
\label{0.7.2.4}
If $A$ is an admissible ring and $\mathfrak J$ is contained in an ideal of
definition of $A$, then $A$ is separated and complete for the
$\mathfrak J$-preadic topology.
\end{proposition}

Indeed, with the notation of (7.2.3), it is immediate that $v\circ u$ is
the identity map of $A$. On the other hand,
$u_n\circ v:\widehat A\to A/\mathfrak J^n$ is the extension by continuity
(for the $\mathfrak J$-preadic topology on $A$ and the discrete topology on
$A/\mathfrak J^n$) of the canonical map $u_n$; in other words, it is
the canonical map of
$\widehat A=\varprojlim_k A/\mathfrak J^k$ onto $A/\mathfrak J^n$;
$u\circ v$ is therefore the projective limit of this sequence of maps,
that is, by definition, the identity map of $\widehat A$; this
proves the proposition.

\begin{corollary}[7.2.5]
\label{0.7.2.5}
Under the hypotheses of (7.2.3), the following conditions are equivalent:
\begin{enumerate}
  \item[$a)$] The homomorphism $u$ is continuous;

\oldpage[0\textsubscript{I}]{64}

  \item[$b)$] The homomorphism $v$ is bicontinuous;
  \item[$c)$] $A$ is a $\mathfrak J$-adic ring.
\end{enumerate}
\end{corollary}

\begin{corollary}[7.2.6]
\label{0.7.2.6}
Let $A$ be an admissible ring and $\mathfrak J$ an ideal of definition of $A$.
For $A$ to be Noetherian, it is necessary and sufficient that $A/\mathfrak J$ be
Noetherian and that $\mathfrak J/\mathfrak J^2$ be an
$(A/\mathfrak J)$-module of finite type.
\end{corollary}

These conditions are obviously necessary. Conversely, suppose they are
satisfied; since by virtue of (7.2.4) $A$ is complete for the
$\mathfrak J$-preadic topology, for it to be Noetherian it is necessary and sufficient that
the associated graded ring $\operatorname{grad}(A)$ (for the filtration by the
$\mathfrak J^n$) be so ([I], p.~18-07, th.~4). Now let
$a_1,\ldots,a_n$ be elements of $\mathfrak J$ whose classes modulo
$\mathfrak J^2$ are generators of $\mathfrak J/\mathfrak J^2$ as an
$A/\mathfrak J$-module. It is immediate by induction that the classes modulo
$\mathfrak J^{m+1}$ of the monomials of total degree $m$ in the $a_i$
$(1\leq i\leq n)$ form a system of generators of the
$(A/\mathfrak J)$-module $\mathfrak J^m/\mathfrak J^{m+1}$. We conclude that
$\operatorname{grad}(A)$ is a ring isomorphic to a quotient of
$(A/\mathfrak J)[T_1,\ldots,T_n]$ ($T_i$ indeterminates), which completes the
proof.

\begin{proposition}[7.2.7]
\label{0.7.2.7}
Let $(A_i,u_{ij})$ be a projective system $(i\in\mathbf N)$ of discrete rings,
and for every integer $i$, let $\mathfrak J_i$ be the kernel in $A_i$ of
the homomorphism $u_{0i}:A_i\to A_0$. Suppose that:
\begin{enumerate}
  \item[$a)$] For $i\leq j$, $u_{ij}$ is surjective and its kernel is
    $\mathfrak J_j^{i+1}$ (so $A_i$ is isomorphic to
    $A_j/\mathfrak J_j^{i+1}$).
  \item[$b)$] $\mathfrak J_1/\mathfrak J_1^2$ ($=\mathfrak J_1$) is a
    module of finite type over $A_0=A_1/\mathfrak J_1$.
\end{enumerate}
Let $A=\varprojlim_i A_i$, and for every integer $n\geq0$, let $u_n$ be
the canonical homomorphism $A\to A_n$ and $\mathfrak J^{(n+1)}\subset A$ its
kernel. Under these conditions:
\begin{enumerate}
  \item[(i)] $A$ is an adic ring, having as an ideal of definition
    $\mathfrak J=\mathfrak J^{(1)}$.
  \item[(ii)] We have $\mathfrak J^{(n)}=\mathfrak J^n$ for every $n\geq1$.
  \item[(iii)] $\mathfrak J/\mathfrak J^2$ is isomorphic to
    $\mathfrak J_1=\mathfrak J_1/\mathfrak J_1^2$, and is consequently a module
    of finite type over $A_0=A/\mathfrak J$.
\end{enumerate}
\end{proposition}

The hypothesis that the $u_{ij}$ are surjective implies that $u_n$ is surjective;
moreover, hypothesis $a)$ implies that $\mathfrak J_j^{j+1}=0$, so $A$ is
an admissible ring (7.2.2); by definition, the $\mathfrak J^{(n)}$ form
a fundamental system of neighbourhoods of $0$ in $A$, so (ii) implies (i).
Moreover, we have $\mathfrak J=\varprojlim_i\mathfrak J_i$ and the maps
$\mathfrak J\to\mathfrak J_i$ are surjective, so (ii) implies (iii), and
we are reduced to proving (ii). By definition, $\mathfrak J^{(n)}$ consists
of the elements $(x_k)_{k\geq0}$ of $A$ such that $x_k=0$ for $k<n$, hence
$\mathfrak J^{(n)}\mathfrak J^{(m)}\subset\mathfrak J^{(n+m)}$; in other words,
the $\mathfrak J^{(n)}$ constitute a \emph{filtration} of $A$. On the other
hand, $\mathfrak J^{(n)}/\mathfrak J^{(n+1)}$ is isomorphic to the projection of
$\mathfrak J^{(n)}$ onto $A_n$; since
$\mathfrak J^{(n)}=\varprojlim_{i>n}\mathfrak J_i^n$, this projection is none
other than $\mathfrak J_n^n$, which is a module over
$A_0=A_n/\mathfrak J_n$. Now let $a_j=(a_{jk})_{k\geq0}$ be $r$ elements of
$\mathfrak J=\mathfrak J^{(1)}$ such that $a_{11},\ldots,a_{r1}$ form a
system of generators of $\mathfrak J_1$ over $A_0$; we shall see that
the set $S_n$ of monomials of total degree $n$ in the $a_j$ generates the ideal
$\mathfrak J^{(n)}$ of $A$. Since $\mathfrak J_i^{i+1}=0$, it is first of all
clear that $S_n\subset\mathfrak J^{(n)}$; since $A$ is complete for the
filtration $(\mathfrak J^{(m)})$, it suffices to prove that the set
$\overline S_n$ of the classes modulo $\mathfrak J^{(n+1)}$ of the elements of $S_n$
generates the graded module $\operatorname{grad}(\mathfrak J^{(n)})$ over
the graded ring $\operatorname{grad}(A)$ for the preceding filtration
([I], p.~18-06, lemma); by virtue of the definition of multiplication in
$\operatorname{grad}(A)$,

\oldpage[0\textsubscript{I}]{65}

it will suffice to prove that for every $m$, $\overline S_m$ is a system of
generators of the $A_0$-module
$\mathfrak J^{(m)}/\mathfrak J^{(m+1)}$, or again that $\mathfrak J_m^m$ is
generated by the monomials of degree $m$ in the $a_{jm}$ ($1\leq j\leq r$).
For this, it remains to show that $\mathfrak J_m$ is (as an
$A_m$-module) generated by the monomials of degree $\leq m$ with respect to the
$a_{jm}$; the proposition being obvious by definition for $m=1$,
let us argue by induction on $m$, and let $\mathfrak J'_m$ be the sub-$A_m$-module
of $\mathfrak J_m$ generated by these monomials. The relation
$\mathfrak J_{m-1}=\mathfrak J_m/\mathfrak J_m^m$ and the induction hypothesis
prove that $\mathfrak J_m=\mathfrak J'_m+\mathfrak J_m^m$, whence, since
$\mathfrak J_m^{m+1}=0$, we obtain
$\mathfrak J_m^m={\mathfrak J'_m}^{\,m}$, and finally
$\mathfrak J_m=\mathfrak J'_m$.

\begin{corollary}[7.2.8]
\label{0.7.2.8}
Under the conditions of (7.2.7), for $A$ to be Noetherian, it is necessary and
sufficient that $A_0$ be so.
\end{corollary}

This follows at once from (7.2.6).

\begin{proposition}[7.2.9]
\label{0.7.2.9}
Suppose the hypotheses of (7.2.7) are satisfied: for every integer $i$, let
$M_i$ be an $A_i$-module, and, for $i\leq j$, let
$v_{ij}:M_j\to M_i$ be a $u_{ij}$-homomorphism, such that $(M_i,v_{ij})$ is a
projective system. Suppose moreover that $M_0$ is an $A_0$-module of finite
type, that $v_{ij}$ is surjective and that its kernel is
$\mathfrak J_j^{i+1}M_j$. Then $M=\varprojlim M_i$ is an $A$-module of finite
type, and the kernel of the surjective $u_n$-homomorphism $v_n:M\to M_n$ is
$\mathfrak J^{n+1}M$ (so that $M_n$ is identified with
$M/\mathfrak J^{n+1}M=M\otimes_A(A/\mathfrak J^{n+1})$).
\end{proposition}

Let $z_h=(z_{hk})_{k\geq0}$ be a system of $s$ elements of $M$ such that the
$z_{h0}$ ($1\leq h\leq s$) form a system of generators of $M_0$; we shall
show that the $z_h$ generate the $A$-module $M$. The $A$-module $M$ is
separated and complete for the filtration formed by the $M^{(n)}$, where $M^{(n)}$
is the set of the $y=(y_k)_{k\geq0}$ in $M$ such that $y_k=0$ for $k<n$;
it is clear that we have $\mathfrak J^{(n)}M\subset M^{(n)}$ and that
$M^{(n)}/M^{(n+1)}=\mathfrak J_n^nM_n$. We are therefore reduced to showing that the
classes of the $z_h$ modulo $M^{(1)}$ generate the graded module
$\operatorname{grad}(M)$ (for the preceding filtration) over the graded ring
$\operatorname{grad}(A)$ ([I], p.~18-06, lemma); for this, one easily
verifies that it again suffices to prove that the $z_{hn}$ ($1\leq h\leq s$)
generate the $A_n$-module $M_n$. We again argue by induction on
$n$, the proposition being obvious by definition for $n=0$; the relation
$M_{n-1}=M_n/\mathfrak J_n^nM_n$ and the induction hypothesis show that if
$M'_n$ is the submodule of $M_n$ generated by the $z_{hn}$, we have
$M_n=M'_n+\mathfrak J_n^nM_n$, and since $\mathfrak J_n$ is nilpotent, this
implies $M_n=M'_n$. The same argument of passing to the associated graded modules
shows that the canonical map from $\mathfrak J^{(n)}M$ to
$M^{(n)}$ is surjective (hence bijective), in other words that
$\mathfrak J^{(n)}M=\mathfrak J^nM$ is the kernel of $M\to M_{n-1}$.

\begin{corollary}[7.2.10]
\label{0.7.2.10}
Let $(N_i,w_{ij})$ be a second projective system of $A_i$-modules satisfying the
conditions of (7.2.9), and let $N=\varprojlim N_i$. There is a one-to-one
correspondence between the projective systems $(h_i)$ of $A_i$-homomorphisms
$h_i:M_i\to N_i$ and the homomorphisms of $A$-modules $h:M\to N$ (which are
necessarily continuous for the $\mathfrak J$-adic topologies).
\end{corollary}

It is clear that if $h:M\to N$ is an $A$-homomorphism, we have
$h(\mathfrak J^nM)\subset\mathfrak J^nN$, whence the continuity of $h$; by
passage to quotients, there therefore corresponds to $h$ a projective system of
$A_i$-homomorphisms $h_i:M_i\to N_i$, of which $h$ is the projective limit,
whence the corollary.

\begin{remark}[7.2.11]
\label{0.7.2.11}
Let $A$ be an adic ring having an ideal of definition $\mathfrak J$ such that
$\mathfrak J/\mathfrak J^2$ is an $(A/\mathfrak J)$-module of finite type;
it is clear that the $A_i=A/\mathfrak J^{i+1}$ satisfy

\oldpage[0\textsubscript{I}]{66}

the conditions of (7.2.7); since $A$ is the projective limit of the $A_i$, we
see that Proposition (7.2.7) gives the description of \emph{all} the adic
rings of the type considered (and in particular of all the \emph{Noetherian
adic} rings).
\end{remark}

\begin{example}[7.2.12]
\label{0.7.2.12}
Let $B$ be a ring and $\mathfrak J$ an ideal of $B$ such that
$\mathfrak J/\mathfrak J^2$ is a module of finite type over
$B/\mathfrak J$ (or over $B$, which amounts to the same); set
$A=\varprojlim_n B/\mathfrak J^{n+1}$; $A$ is the separated completion of $B$
endowed with the $\mathfrak J$-preadic topology. If
$A_n=B/\mathfrak J^{n+1}$, it is immediate that the $A_n$ satisfy the
conditions of (7.2.7); hence $A$ is an adic ring, and if
$\overline{\mathfrak J}$ is the closure in $A$ of the canonical image of
$\mathfrak J$, then $\overline{\mathfrak J}$ is an ideal of definition of $A$,
$\overline{\mathfrak J^n}$ is the closure of the canonical image of
$\mathfrak J^n$, $A/\overline{\mathfrak J^n}$ is identified with
$B/\mathfrak J^n$, and
$\overline{\mathfrak J}/\overline{\mathfrak J^2}$ is isomorphic to
$\mathfrak J/\mathfrak J^2$ as an
$(A/\overline{\mathfrak J})$-module. Likewise, if $N$ is such that
$N/\mathfrak JN$ is a $B$-module of finite type, and if we set
$M_i=N/\mathfrak J^{i+1}N$, then $M=\varprojlim M_i$ is an $A$-module of finite
type, isomorphic to the separated completion of $N$ for the
$\mathfrak J$-preadic topology, $\overline{\mathfrak J^n}M$ is identified with
the closure of the canonical image of $\mathfrak J^nN$, and
$M/\overline{\mathfrak J^n}M$ with $N/\mathfrak J^nN$.
\end{example}

\subsection{Noetherian preadic rings}
\label{subsection:0.7.3}

\begin{env}[7.3.1]
\label{0.7.3.1}
Let $A$ be a ring, $\mathfrak J$ an ideal of $A$, and $M$ an $A$-module;
we shall denote by
$\widehat A=\varprojlim_n A/\mathfrak J^n$ (resp.
$\widehat M=\varprojlim_n M/\mathfrak J^nM$) the separated completion of $A$
(resp. $M$) for the $\mathfrak J$-preadic topology. Let
\[
  M'\xrightarrow{u}M\xrightarrow{v}M''\to0
\]
be an exact sequence of $A$-modules; since
$M/\mathfrak J^nM=M\otimes_A(A/\mathfrak J^n)$, the sequence
\[
  M'/\mathfrak J^nM'\xrightarrow{u_n}
  M/\mathfrak J^nM\xrightarrow{v_n}
  M''/\mathfrak J^nM''\to0
\]
is exact for every $n$. Moreover, since
$v(\mathfrak J^nM)=\mathfrak J^nv(M)=\mathfrak J^nM''$,
$\widehat v=\varprojlim v_n$ is surjective (Bourbaki,
\emph{Top. gén.}, chap.~IX, 2nd ed., p.~60, cor.~2).
On the other hand, if $z=(z_k)$ is an element of the kernel of $\widehat v$, then for every
integer $k$ there exists a $z'_k\in M'/\mathfrak J^kM'$ such that
$u_k(z'_k)=z_k$; we conclude that there exists
$z'=(z'_n)\in\widehat{M'}$ such that the first $k$ components of
$\widehat u(z')$ coincide with those of $z$; in other words, the image under
$\widehat u$ of $\widehat{M'}$ is \emph{dense} in the kernel of
$\widehat v$.

If $A$ is assumed \emph{Noetherian}, so is $\widehat A$, by
virtue of (7.2.12), since $\mathfrak J/\mathfrak J^2$ is then an $A$-module of
finite type. Moreover:
\end{env}

\begin{namedtheorem}[7.3.2]{Krull's theorem}
\label{0.7.3.2}
Let $A$ be a Noetherian ring, $\mathfrak J$ an ideal of $A$, $M$ an
$A$-module of finite type, and $M'$ a submodule of $M$; then the topology
induced on $M'$ by the $\mathfrak J$-preadic topology of $M$ is identical
to the $\mathfrak J$-preadic topology of $M'$.
\end{namedtheorem}

This follows at once from the

\begin{namedlemma}[7.3.2.1]{Artin--Rees lemma}
\label{0.7.3.2.1}
Under the hypotheses of (7.3.2), there exists an integer $p$ such that, for
$n\geq p$, we have
\[
  M'\cap\mathfrak J^nM=\mathfrak J^{n-p}
  (M'\cap\mathfrak J^pM).
\]
\end{namedlemma}

For the proof, see ([I], p.~2-04).

\oldpage[0\textsubscript{I}]{67}

\begin{corollary}[7.3.3]
\label{0.7.3.3}
Under the hypotheses of (7.3.2), the canonical map
$M\otimes_A\widehat A\to\widehat M$ is bijective, and the functor
$M\otimes_A\widehat A$ is exact in $M$ in the category of $A$-modules of
finite type; consequently, the $\mathfrak J$-adic separated completion $\widehat A$
is a flat $A$-module (6.1.1).
\end{corollary}

Let us first note that in the category of $A$-modules of finite type,
$\widehat M$ is an \emph{exact} functor in $M$. Indeed, let
\[
  0\to M'\xrightarrow{u}M\xrightarrow{v}M''\to0
\]
be an exact sequence; we already know that
$\widehat v:\widehat M\to\widehat{M''}$ is surjective (7.3.1); on the other hand,
if $i$ is the canonical homomorphism $M\to\widehat M$, it follows from Krull's
theorem that the closure in $\widehat M$ of $i(u(M'))$ is identified with the separated
completion of $M'$ for the $\mathfrak J$-preadic topology; hence
$\widehat u$ is injective, and by virtue of (7.3.1) the image of $\widehat u$ is
equal to the kernel of $\widehat v$.

That being so, the canonical map $M\otimes_A\widehat A\to\widehat M$
is obtained by passing to the projective limit over the maps
\[
  M\otimes_A\widehat A\to
  M\otimes_A(A/\mathfrak J^n)=M/\mathfrak J^nM.
\]
It is clear that this map is bijective when $M$ is of the form
$A^p$. If $M$ is an $A$-module of finite type, we have an exact sequence
$A^p\to A^q\to M\to0$, whence, by virtue of the \emph{right} exactness of the
functors $M\otimes_A\widehat A$ and $\widehat M$ (in $M$) in the category
of $A$-modules of finite type, the commutative diagram
\[
\begin{tikzcd}[column sep=large,row sep=large]
  A^p\otimes_A\widehat A \arrow[r] \arrow[d] &
  A^q\otimes_A\widehat A \arrow[r] \arrow[d] &
  M\otimes_A\widehat A \arrow[r] \arrow[d] &
  0 \\
  \widehat A^{\,p} \arrow[r] &
  \widehat A^{\,q} \arrow[r] &
  \widehat M \arrow[r] &
  0
\end{tikzcd}
\]
where the two rows are exact and the first two vertical arrows are
isomorphisms; the conclusion follows at once.

\begin{corollary}[7.3.4]
\label{0.7.3.4}
Let $A$ be a Noetherian ring, $\mathfrak J$ an ideal of $A$, and $M$, $N$ two
$A$-modules of finite type; we have canonical functorial isomorphisms
\[
  (M\otimes_A N)^\wedge
  \xrightarrow{\sim}
  \widehat M\otimes_{\widehat A}\widehat N,
  \qquad
  \bigl(\operatorname{Hom}_A(M,N)\bigr)^\wedge
  \xrightarrow{\sim}
  \operatorname{Hom}_{\widehat A}(\widehat M,\widehat N).
\]
\end{corollary}

This follows from (7.3.3), (6.2.1) and (6.2.2).

\begin{corollary}[7.3.5]
\label{0.7.3.5}
Let $A$ be a Noetherian ring and $\mathfrak J$ an ideal of $A$. The
following conditions are equivalent:
\begin{enumerate}
  \item[$a)$] $\mathfrak J$ is contained in the radical of $A$.
  \item[$b)$] $\widehat A$ is a faithfully flat $A$-module (6.4.1).
  \item[$c)$] Every $A$-module of finite type is separated for the
  $\mathfrak J$-preadic topology.
  \item[$d)$] Every submodule of an $A$-module of finite type is closed for the
  $\mathfrak J$-preadic topology.
\end{enumerate}
\end{corollary}

Since $\widehat A$ is a flat $A$-module, conditions $b)$ and $c)$ are
equivalent, because $b)$ amounts to saying that if $M$ is an $A$-module of finite
type, the canonical map
$M\to\widehat M=M\otimes_A\widehat A$ is injective (6.6.1, $c)$). It is
immediate that $c)$ implies $d)$, because if $N$ is a submodule of an $A$-module
$M$ of finite type, $M/N$ is separated for the $\mathfrak J$-preadic topology,
so $N$ is closed in $M$. Let us show that $d)$ implies $a)$: if
$\mathfrak m$ is a maximal ideal of $A$, $\mathfrak m$ is closed in $A$
for the $\mathfrak J$-preadic topology, so
\[
  \mathfrak m=\bigcap_{p\geq0}(\mathfrak m+\mathfrak J^p),
\]
and since $\mathfrak m+\mathfrak J^p$ is necessarily equal to $A$ or to
$\mathfrak m$, we have $\mathfrak m+\mathfrak J^p=\mathfrak m$ for $p$ sufficiently

\oldpage[0\textsubscript{I}]{68}

large, whence $\mathfrak J^p\subset\mathfrak m$ and
$\mathfrak J\subset\mathfrak m$ since $\mathfrak m$ is prime. Finally,
$a)$ implies $b)$: indeed, let $P$ be the closure of $\{0\}$ in an
$A$-module $M$ of finite type, for the $\mathfrak J$-preadic topology; by
virtue of Krull's theorem (7.3.2), the topology induced on $P$ by the
$\mathfrak J$-preadic topology of $M$ is the $\mathfrak J$-preadic topology of
$P$, so $\mathfrak JP=P$; since $P$ is of finite type, it follows from Nakayama's
lemma that $P=0$ ($\mathfrak J$ being contained in the radical of $A$).

Note that the conditions of (7.3.5) are satisfied when $A$ is a
\emph{Noetherian local ring} and $\mathfrak J\neq A$ is an arbitrary ideal of
$A$.

\begin{corollary}[7.3.6]
\label{0.7.3.6}
If $A$ is a Noetherian $\mathfrak J$-adic ring, every $A$-module of finite
type is separated and complete for the $\mathfrak J$-preadic topology.
\end{corollary}

Since we then have $\widehat A=A$, this follows at once from (7.3.3).

We conclude that Proposition (7.2.9) gives the description of \emph{all} the
modules of finite type over a Noetherian adic ring.

\begin{corollary}[7.3.7]
\label{0.7.3.7}
Under the hypotheses of (7.3.2), the kernel of the canonical map
$M\to\widehat M=M\otimes_A\widehat A$ is the set of $x\in M$ annihilated by
an element of $1+\mathfrak J$.
\end{corollary}

Indeed, for $x\in M$ to belong to this kernel, it is necessary and sufficient that
the separated completion of the submodule $Ax$ reduce to 0 (by (7.3.2)), in other
words that $x\in\mathfrak Jx$.

\subsection{Quasi-finite modules over local rings}
\label{subsection:0.7.4}

\begin{definition}[7.4.1]
\label{0.7.4.1}
Given a local ring $A$ with maximal ideal $\mathfrak m$, an
$A$-module $M$ is said to be quasi-finite (over $A$) if $M/\mathfrak mM$ is of finite rank
over the residue field $k=A/\mathfrak m$.
\end{definition}

When $A$ is Noetherian, the separated completion $\widehat M$ of $M$ for the
$\mathfrak m$-preadic topology is then an
$\widehat A$-\emph{module of finite type}; indeed, since
$\mathfrak m/\mathfrak m^2$ is then an $A$-module of finite type, this follows
from (7.2.12) and from the hypothesis on $M/\mathfrak mM$.

In particular, if we suppose moreover that $A$ is \emph{complete} and $M$
\emph{separated} for the $\mathfrak m$-preadic topology (in other words,
$\bigcap_n\mathfrak m^nM=0$), then $M$ itself is an
$A$-\emph{module of finite type}: indeed, $\widehat M$ is then an
$A$-module of finite type, and since $M$ is identified with a submodule of
$\widehat M$, $M$ is also of finite type (and moreover identical to its
completion by virtue of (7.3.6)).

\begin{proposition}[7.4.2]
\label{0.7.4.2}
Let $A$, $B$ be two local rings, $\mathfrak m$, $\mathfrak n$ their maximal
ideals, and suppose $B$ Noetherian. Let $\varphi:A\to B$ be a
local homomorphism and $M$ a $B$-module of finite type. If $M$ is a quasi-finite
$A$-module, the $\mathfrak m$-preadic and $\mathfrak n$-preadic topologies
on $M$ are identical, hence separated.
\end{proposition}

Let us observe that by hypothesis $M/\mathfrak mM$ is of \emph{finite length} as
an $A$-module, hence also \emph{a fortiori} as a $B$-module. We
conclude that $\mathfrak n$ is the \emph{only prime ideal} of $B$ containing
the annihilator of $M/\mathfrak mM$: indeed, we immediately reduce, by virtue
of (I.7.4) and (I.7.2), to the case where $M/\mathfrak mM$ is \emph{simple}, hence
necessarily isomorphic to $B/\mathfrak n$, and our assertion is obvious
in this last case. On the other hand, since $M$ is a $B$-module of finite type,
the prime ideals that contain the annihilator of $M/\mathfrak mM$ are those
that contain $\mathfrak mB+\mathfrak b$, where $\mathfrak b$ denotes
the annihilator of the $B$-module $M$ (I.7.5). Since $B$ is Noetherian, we conclude
([III], p.~127, cor.~4) that $\mathfrak mB+\mathfrak b$ is an ideal

\oldpage[0\textsubscript{I}]{69}

of definition of $B$, in other words that there exists $k>0$ such that
$\mathfrak n^k\subset\mathfrak mB+\mathfrak b\subset\mathfrak n$; consequently,
for every $h>0$
\[
  \mathfrak n^{hk}M
  \subset(\mathfrak mB+\mathfrak b)^hM
  =\mathfrak m^hM
  \subset\mathfrak n^hM,
\]
which proves that the $\mathfrak m$-preadic and
$\mathfrak n$-preadic topologies on $M$ are the same; the second is moreover
separated by virtue of (7.3.5).

\begin{corollary}[7.4.3]
\label{0.7.4.3}
Under the hypotheses of (7.4.2), if moreover $A$ is Noetherian and complete for
the $\mathfrak m$-preadic topology, $M$ is an $A$-module of finite type.
\end{corollary}

Indeed, $M$ is then separated for the $\mathfrak m$-preadic topology, and
our assertion follows from the remark made in (7.4.1).

\begin{env}[7.4.4]
\label{0.7.4.4}
The most important case of application of (7.4.2) is the one where $B$
itself is a \emph{quasi-finite $A$-module}, which amounts to saying that
$B/\mathfrak mB$ is an \emph{algebra of finite rank over}
$k=A/\mathfrak m$; this condition can moreover be decomposed into the
conjunction of the following two, by virtue of what precedes:
\begin{enumerate}
  \item[(i)] $\mathfrak mB$ \emph{is an ideal of definition of} $B$;
  \item[(ii)] $B/\mathfrak n$ \emph{is an extension of finite rank of the field}
  $A/\mathfrak m$.
\end{enumerate}
When this is so, every $B$-module of finite type is obviously a
quasi-finite $A$-module.
\end{env}

\begin{corollary}[7.4.5]
\label{0.7.4.5}
Under the hypotheses of (7.4.2), if $\mathfrak b$ is the annihilator of the
$B$-module $M$, $B/\mathfrak b$ is a quasi-finite $A$-module.
\end{corollary}

Suppose $M\neq0$ (otherwise the corollary is obvious). We can consider $M$
as a module over the Noetherian local ring $B/\mathfrak b$; its annihilator
then being reduced to $0$, the proof of (7.4.2) shows that
$\mathfrak m(B/\mathfrak b)$ is an ideal of definition of
$B/\mathfrak b$. Furthermore, $M/\mathfrak nM$ is a vector space of
finite rank over $A/\mathfrak m$, being a quotient of $M/\mathfrak mM$, which is
by hypothesis of finite rank over $A/\mathfrak m$; since $M\neq0$, we have
$M\neq\mathfrak nM$ by virtue of Nakayama's lemma; since
$M/\mathfrak nM$ is a vector space $\neq0$ over $B/\mathfrak n$, the fact
that it is of finite rank over $A/\mathfrak m$ implies that $B/\mathfrak n$ is
also of finite rank over $A/\mathfrak m$; the conclusion therefore follows from (7.4.4)
applied to the ring $B/\mathfrak b$.


\subsection{Rings of restricted formal power series}
\label{subsection:0.7.5}

\begin{env}[7.5.1]
\label{0.7.5.1}
Let $A$ be a topological ring that is linearly topologized, separated and complete;
let $(\mathfrak J_\lambda)$ be a fundamental system of neighbourhoods of $0$ in
$A$ consisting of (open) ideals, so that $A$ is canonically identified with
$\varprojlim A/\mathfrak J_\lambda$ (7.2.1). For every $\lambda$, let
$B_\lambda=(A/\mathfrak J_\lambda)[T_1,\ldots,T_r]$, where the $T_i$ are
indeterminates; it is clear that the $B_\lambda$ form a projective system
of discrete rings. We shall set
$\varprojlim B_\lambda=A\{T_1,\ldots,T_r\}$, and we shall see that this
topological ring is independent of the fundamental system of ideals
$(\mathfrak J_\lambda)$ considered. More precisely, let $A'$ be the subring
of the ring of formal power series $A[\![T_1,\ldots,T_r]\!]$ consisting of the formal
power series $\sum_\alpha c_\alpha T^\alpha$ (with
$\alpha=(\alpha_1,\ldots,\alpha_r)\in\mathbf{N}^r$) such that
$\lim_\alpha c_\alpha=0$ (along the filter of complements of finite
subsets of $\mathbf{N}^r$); we shall say that these series are the
\emph{restricted} formal power series in the $T_i$, with coefficients in $A$.

\oldpage[0\textsubscript{I}]{70}

For every neighbourhood $V$ of $0$ in $A$, let $V'$ be the set of
$x=\sum_\alpha c_\alpha T^\alpha\in A'$ such that $c_\alpha\in V$ for every
$\alpha$. One verifies at once that the $V'$ form a fundamental system of
neighbourhoods of $0$ defining on $A'$ a \emph{separated} ring topology;
we shall define canonically a \emph{topological isomorphism} of
the ring $A\{T_1,\ldots,T_r\}$ onto $A'$. For every
$\alpha\in\mathbf{N}^r$ and every $\lambda$, let
$\varphi_{\lambda,\alpha}$ be the map from
$(A/\mathfrak J_\lambda)[T_1,\ldots,T_r]$ to
$A/\mathfrak J_\lambda$ which assigns to every polynomial of the first ring
the coefficient of $T^\alpha$ in this polynomial. It is clear that
the $\varphi_{\lambda,\alpha}$ form a projective system of homomorphisms
of $(A/\mathfrak J_\lambda)$-modules, whose projective limit is a
continuous homomorphism
$\varphi_\alpha:A\{T_1,\ldots,T_r\}\to A$; we shall see that, for every
$y\in A\{T_1,\ldots,T_r\}$, the formal power series
$\sum_\alpha\varphi_\alpha(y)T^\alpha$ is \emph{restricted}. Indeed, if
$y_\lambda$ is the component of $y$ in $B_\lambda$, and if we denote by
$H_\lambda$ the finite set of the $\alpha\in\mathbf{N}^r$ for which the
coefficients of the polynomial $y_\lambda$ are not zero, we have
$\varphi_{\mu,\alpha}(y_\mu)\in\mathfrak J_\lambda$ for
$\mathfrak J_\mu\subset\mathfrak J_\lambda$ and $\alpha\notin H_\lambda$, and,
by passage to the limit, $\varphi_\alpha(y)\in\mathfrak J_\lambda$ for
$\alpha\notin H_\lambda$. We therefore define a ring homomorphism
$\varphi:A\{T_1,\ldots,T_r\}\to A'$ by setting
$\varphi(y)=\sum_\alpha\varphi_\alpha(y)T^\alpha$, and it is immediate that
$\varphi$ is continuous. Conversely, if $\theta_\lambda$ is the canonical
homomorphism $A\to A/\mathfrak J_\lambda$, then for every element
$z=\sum_\alpha c_\alpha T^\alpha\in A'$ and every $\lambda$, there is only a
finite number of indices $\alpha$ such that $\theta_\lambda(c_\alpha)\neq0$, and
consequently
$\psi_\lambda(z)=\sum_\alpha\theta_\lambda(c_\alpha)T^\alpha$ belongs to
$B_\lambda$; the $\psi_\lambda$ are continuous and form a projective
system of homomorphisms whose projective limit is a continuous
homomorphism $\psi:A'\to A\{T_1,\ldots,T_r\}$; finally, it remains to verify that
$\varphi\circ\psi$ and $\psi\circ\varphi$ are the identity automorphisms,
which is immediate.
\end{env}

\begin{env}[7.5.2]
\label{0.7.5.2}
We shall identify $A\{T_1,\ldots,T_r\}$ with the ring $A'$ of restricted formal
power series by means of the isomorphisms defined in (7.5.1). The canonical
isomorphisms
\[
  ((A/\mathfrak J_\lambda)[T_1,\ldots,T_r])
  [T_{r+1},\ldots,T_s]
  \xrightarrow{\sim}
  (A/\mathfrak J_\lambda)[T_1,\ldots,T_s]
\]
define, by passage to the projective limit, a canonical isomorphism
\[
  (A\{T_1,\ldots,T_r\})\{T_{r+1},\ldots,T_s\}
  \xrightarrow{\sim}
  A\{T_1,\ldots,T_s\}
\]
\end{env}

\begin{env}[7.5.3]
\label{0.7.5.3}
For every continuous homomorphism $u:A\to B$ from $A$ to a linearly
topologized ring $B$ that is separated and complete, and every system $(b_1,\ldots,b_r)$ of $r$
elements of $B$, there exists \emph{one and only one continuous homomorphism}
$\overline u:A\{T_1,\ldots,T_r\}\to B$ such that $\overline u(a)=u(a)$ for
every $a\in A$ and $\overline u(T_j)=b_j$ for $1\leq j\leq r$. Indeed, it suffices
to take
\[
  \overline u\Bigl(\sum_\alpha c_\alpha T^\alpha\Bigr)
  =\sum_\alpha u(c_\alpha)b_1^{\alpha_1}\cdots b_r^{\alpha_r}~;
\]
the verifications of the fact that the family
$(u(c_\alpha)b_1^{\alpha_1}\cdots b_r^{\alpha_r})$ is summable in $B$ and
of the continuity of $\overline u$ are immediate and are left to the reader. Note
that this property (for arbitrary $B$ and $b_j$)
\emph{characterizes} the topological ring $A\{T_1,\ldots,T_r\}$ up to
unique isomorphism.
\end{env}

\begin{proposition}[7.5.4]
\label{0.7.5.4}
\begin{enumerate}
  \item[(i)] If $A$ is an admissible ring, so is
  $A'=A\{T_1,\ldots,T_r\}$.
  \item[(ii)] Let $A$ be an adic ring and $\mathfrak J$ an ideal of definition
  of $A$ such that $\mathfrak J/\mathfrak J^2$ is of finite type

\oldpage[0\textsubscript{I}]{71}

  over $A/\mathfrak J$. If we set $\mathfrak J'=\mathfrak JA'$, then $A'$ is
  a $\mathfrak J'$-adic ring, and
  $\mathfrak J'/\mathfrak J'^{\,2}$ is of finite type over
  $A'/\mathfrak J'$. If, moreover, $A$ is Noetherian, so is
  $A'$.
\end{enumerate}
\end{proposition}

\begin{enumerate}
  \item[(i)] If $\mathfrak J$ is an ideal of $A$ and $\mathfrak J'$ the ideal of
  $A'$ consisting of the $\sum_\alpha c_\alpha T^\alpha$ such that
  $c_\alpha\in\mathfrak J$ for every $\alpha$, then
  $(\mathfrak J')^n\subset(\mathfrak J^n)'$; if $\mathfrak J$ is an ideal of
  definition of $A$, $\mathfrak J'$ is therefore an ideal of definition of $A'$.

  \item[(ii)] Set $A_i=A/\mathfrak J^{i+1}$, and for $i\leq j$, let
  $u_{ij}$ be the canonical homomorphism
  $A/\mathfrak J^{j+1}\to A/\mathfrak J^{i+1}$; set
  $A'_i=A_i[T_1,\ldots,T_r]$, and let $u'_{ij}$ be the homomorphism
  $A'_j\to A'_i$ ($i\leq j$) obtained by applying $u_{ij}$ to the coefficients
  of the polynomials of $A'_j$. Let us show that the projective system
  $(A'_i,u'_{ij})$ satisfies the conditions of (7.2.7); since $\mathfrak J'$ is
  the kernel of $A'\to A'_0$, this will prove the first assertion of (ii). Now,
  it is clear that the $u'_{ij}$ are surjective; the kernel $\mathfrak J'_i$
  of $u'_{0i}$ is the set of polynomials of
  $A_i[T_1,\ldots,T_r]$ whose coefficients lie in
  $\mathfrak J/\mathfrak J^{i+1}$; in particular, $\mathfrak J'_1$ is
  the set of polynomials of $A_1[T_1,\ldots,T_r]$ whose coefficients
  lie in $\mathfrak J/\mathfrak J^2$. Since $\mathfrak J/\mathfrak J^2$
  is of finite type over $A_1=A/\mathfrak J^2$, we see already that
  $\mathfrak J'_1/\mathfrak J_1'^{\,2}$ is a module of finite type over $A'_1$
  (or, what amounts to the same, over
  $A'_0=A'_1/\mathfrak J'_1$). Let us show that the kernel of $u'_{ij}$ is
  $\mathfrak J_j'^{\,i+1}$. It is obvious that
  $\mathfrak J_j'^{\,i+1}$ is contained in this kernel. On the other hand, let
  $a_1,\ldots,a_m$ be elements of $\mathfrak J$ whose classes
  modulo~$\mathfrak J^2$ generate $\mathfrak J/\mathfrak J^2$; one verifies
  immediately that the classes modulo~$\mathfrak J^{j+1}$ of the monomials of degree
  $\leq j$ in the $a_k$ ($1\leq k\leq m$) generate
  $\mathfrak J/\mathfrak J^{j+1}$, and the classes of the monomials of degree $>i$
  and $\leq j$ therefore generate
  $\mathfrak J^{i+1}/\mathfrak J^{j+1}$; a monomial in the $T_k$ having such an
  element as coefficient is therefore a product of $i+1$ elements of
  $\mathfrak J'_j$, which establishes our assertion. Finally, if $A$ is
  Noetherian, so is
  $A'/\mathfrak J'=(A/\mathfrak J)[T_1,\ldots,T_r]$, hence $A'$ is Noetherian
  (7.2.8).
\end{enumerate}

\begin{proposition}[7.5.5]
\label{0.7.5.5}
Let $A$ be a Noetherian $\mathfrak J$-adic ring, $B$ an admissible topological
ring, and $\varphi:A\to B$ a continuous homomorphism making $B$ an
$A$-algebra. The following conditions are equivalent:
\begin{enumerate}
  \item[a)] $B$ is Noetherian and $\mathfrak JB$-adic, and
  $B/\mathfrak JB$ is an algebra of finite type over $A/\mathfrak J$.

  \item[b)] $B$ is topologically $A$-isomorphic to $\varprojlim B_n$, where
  $B_n=B_m/\mathfrak J^{n+1}B_m$ for $m\geq n$, and $B_1$ is an algebra of
  finite type over $A_1=A/\mathfrak J^2$.

  \item[c)] $B$ is topologically $A$-isomorphic to a quotient of an algebra
  of the form\linebreak $A\{T_1,\ldots,T_r\}$ by an ideal (necessarily closed by virtue
  of (7.3.6) and (7.5.4 (ii))).
\end{enumerate}
\end{proposition}

Since $A$ is Noetherian, so is
$A'=A\{T_1,\ldots,T_r\}$ (7.5.4), hence c) implies that $B$ is Noetherian;
since $\mathfrak J'=\mathfrak JA'$ is an open neighbourhood of 0 in $A'$ such
that the $\mathfrak J'^{n}$ form a fundamental system of neighbourhoods of 0,
the images $\mathfrak J^nB$ of the $\mathfrak J'^{n}$ form a fundamental
system of neighbourhoods of 0 in $B$, and since we know that $B$ is separated and
complete, $B$ is a $\mathfrak JB$-adic ring. Finally, $B/\mathfrak JB$ is
an algebra (over $A/\mathfrak J$) that is a quotient of
$A'/\mathfrak JA'=(A/\mathfrak J)[T_1,\ldots,T_r]$, hence is of finite type, which
completes the proof that c) implies a).

If $B$ is $\mathfrak JB$-adic and Noetherian, $B$ is isomorphic to
$\varprojlim B_n$, where $B_n=B/\mathfrak J^{n+1}B$ (7.2.11), and
$\mathfrak JB/\mathfrak J^2B$ is a module of finite type over
$B/\mathfrak JB$. Let $(a_j)_{1\leq j\leq s}$ be a system of generators of the
$(B/\mathfrak JB)$-module $\mathfrak JB/\mathfrak J^2B$, and
$(c_i)_{1\leq i\leq r}$ a system of elements of $B/\mathfrak J^2B$ whose
classes

\oldpage[0\textsubscript{I}]{72}

modulo~$\mathfrak JB/\mathfrak J^2B$ form a system of generators of the
$(A/\mathfrak J)$-algebra $B/\mathfrak JB$; one sees at once that the $c_i,a_j$
form a system of generators of the $(A/\mathfrak J^2)$-algebra
$B/\mathfrak J^2B$, hence a) implies b).

It remains to prove that b) implies c). The hypothesis implies that $B_1$ is a
Noetherian ring, and since $B_1=B_2/\mathfrak J^2B_2$, we have
$\mathfrak J^2B_1=0$, so
$\mathfrak JB_1=\mathfrak JB_1/\mathfrak J^2B_1$ is a $B_0$-module of finite
type. The conditions of (7.2.7) are therefore satisfied by the projective system
$(B_n)$, and $B$ is a $\mathfrak JB$-adic ring. Let
$(c_i)_{1\leq i\leq r}$ be a finite system of elements of $B$ whose classes
modulo~$\mathfrak JB$ generate the $(A/\mathfrak J)$-algebra
$B/\mathfrak JB$, and whose linear combinations with coefficients in
$\mathfrak J$ are such that their classes modulo~$\mathfrak J^2B$ generate
the $B_0$-module $\mathfrak JB/\mathfrak J^2B$. There exists a continuous
$A$-homomorphism $u$ from $A'=A\{T_1,\ldots,T_r\}$ to $B$ which
reduces to $\varphi$ on $A$ and is such that $u(T_i)=c_i$ for $1\leq i\leq r$
(7.5.3); if we prove that $u$ is \emph{surjective}, c) will be established, for from
$u(A')=B$ we shall deduce $u(\mathfrak J^nA')=\mathfrak J^nB$; in other words, $u$
will be a strict morphism of topological rings, and $B$ will therefore be isomorphic to a
quotient of $A'$ by a closed ideal. Now, since $B$ is complete for the
$\mathfrak JB$-adic topology, it suffices ([I], p.~18-07) to show that
the homomorphism $\operatorname{grad}(A')\to\operatorname{grad}(B)$ deduced
canonically from $u$ for the $\mathfrak J$-adic filtrations on $A'$ and
$B$ is surjective. But by definition, the homomorphisms
$A'/\mathfrak JA'\to B/\mathfrak JB$ and
$\mathfrak JA'/\mathfrak J^2A'\to\mathfrak JB/\mathfrak J^2B$ deduced from $u$
are surjective; by induction on $n$, one deduces at once that the same holds for
$\mathfrak JA'/\mathfrak J^nA'\to\mathfrak JB/\mathfrak J^nB$, and
\emph{a fortiori} for
$\mathfrak J^nA'/\mathfrak J^{n+1}A'\to
\mathfrak J^nB/\mathfrak J^{n+1}B$, which completes the proof.

\subsection{Complete rings of fractions}
\label{subsection:0.7.6}

\begin{env}[7.6.1]
\label{0.7.6.1}
Let $A$ be a linearly topologized ring, $(\mathfrak J_\lambda)$ a
fundamental system of neighbourhoods of 0 in $A$ consisting of ideals, and $S$ a
multiplicative subset of $A$. Let $u_\lambda$ be the canonical homomorphism
$A\to A_\lambda=A/\mathfrak J_\lambda$, and for
$\mathfrak J_\mu\subset\mathfrak J_\lambda$, let $u_{\lambda\mu}$ be
the canonical homomorphism $A_\mu\to A_\lambda$. Set
$S_\lambda=u_\lambda(S)$, so that
$u_{\lambda\mu}(S_\mu)=S_\lambda$. The $u_{\lambda\mu}$ canonically give
surjective homomorphisms
$S_\mu^{-1}A_\mu\to S_\lambda^{-1}A_\lambda$, for which these rings
form a projective system; denote by $A\{S^{-1}\}$ the projective
limit of this system. This definition does not depend on the fundamental system
of neighbourhoods $(\mathfrak J_\lambda)$ chosen; indeed:
\end{env}

\begin{proposition}[7.6.2]
\label{0.7.6.2}
The ring $A\{S^{-1}\}$ is topologically isomorphic to the separated completion of
the ring $S^{-1}A$ for the topology of which a fundamental system of neighbourhoods
of 0 is formed by the $S^{-1}\mathfrak J_\lambda$.
\end{proposition}

Indeed, if $v_\lambda$ is the canonical homomorphism
$S^{-1}A\to S_\lambda^{-1}A_\lambda$ deduced from $u_\lambda$, the kernel of
$v_\lambda$ is $S^{-1}\mathfrak J_\lambda$ and $v_\lambda$ is surjective,
whence the proposition (7.2.1).

\begin{corollary}[7.6.3]
\label{0.7.6.3}
If $S'$ is the canonical image of $S$ in the separated completion $\widehat A$ of
$A$, then $A\{S^{-1}\}$ is canonically identified with
$\widehat A\{{S'}^{-1}\}$.
\end{corollary}

Note that even if $A$ is separated and complete, the same is not true of
$S^{-1}A$ for the topology defined by the
$S^{-1}\mathfrak J_\lambda$, as one sees for example by taking for $S$
the set of the $f^n$ ($n\geq0$), where $f$ is topologically nilpotent and not
nilpotent: indeed, $S^{-1}A$ does not reduce to 0, and on the other hand, for every
$\lambda$ there exists $n$ such that $f^n\in\mathfrak J_\lambda$, hence
$1=f^n/f^n\in S^{-1}\mathfrak J_\lambda$ and
$S^{-1}\mathfrak J_\lambda=S^{-1}A$.

\oldpage[0\textsubscript{I}]{73}

\begin{corollary}[7.6.4]
\label{0.7.6.4}
If, in $A$, $0$ is not adherent to $S$, the ring $A\{S^{-1}\}$ does not
reduce to $0$.
\end{corollary}

Indeed, $0$ is not adherent to $\{1\}$ in the ring $S^{-1}A$; otherwise,
we would have $1\in S^{-1}\mathfrak J_\lambda$ for every open ideal
$\mathfrak J_\lambda$ of $A$, and it would follow that
$\mathfrak J_\lambda\cap S\neq\varnothing$ for every $\lambda$, contrary
to the hypothesis.

\begin{env}[7.6.5]
\label{0.7.6.5}
We shall say that $A\{S^{-1}\}$ is the \emph{complete ring of fractions} of $A$
with denominators in $S$. With the preceding notation, it is clear
that the inverse image of $S^{-1}\mathfrak J_\lambda$ in $A$ contains
$\mathfrak J_\lambda$, so the canonical map $A\to S^{-1}A$ is
continuous, and if it is composed with the canonical map
$S^{-1}A\to A\{S^{-1}\}$, one obtains a continuous canonical homomorphism
$A\to A\{S^{-1}\}$, the projective limit of the homomorphisms
$A\to S_\lambda^{-1}A_\lambda$.
\end{env}

\begin{env}[7.6.6]
\label{0.7.6.6}
The pair formed by $A\{S^{-1}\}$ and the canonical map
$A\to A\{S^{-1}\}$ is characterized by the following \emph{universal
property}: every continuous homomorphism $u$ from $A$ to a linearly
topologized ring $B$ that is separated and complete, such that $u(S)$ consists of elements
invertible in $B$, factors in one and only one way as
$A\to A\{S^{-1}\}\xrightarrow{u'}B$, where $u'$ is continuous. Indeed, $u$
factors in one and only one way as $A\to S^{-1}A\xrightarrow{v'}B$; since for
every open ideal $\mathfrak K$ of $B$, $u^{-1}(\mathfrak K)$ contains some
$\mathfrak J_\lambda$, ${v'}^{-1}(\mathfrak K)$ necessarily contains
$S^{-1}\mathfrak J_\lambda$, so $v'$ is continuous; since $B$ is separated and
complete, $v'$ factors in one and only one way as
$S^{-1}A\to A\{S^{-1}\}\xrightarrow{u'}B$, where $u'$ is continuous; whence our
assertion.
\end{env}

\begin{env}[7.6.7]
\label{0.7.6.7}
Let $B$ be a second linearly topologized ring, $T$ a multiplicative
subset of $B$, and $\varphi:A\to B$ a continuous homomorphism such that
$\varphi(S)\subset T$. By what precedes, the continuous homomorphism
$A\xrightarrow{\varphi}B\to B\{T^{-1}\}$ factors uniquely as
$A\to A\{S^{-1}\}\xrightarrow{\varphi'}B\{T^{-1}\}$, where $\varphi'$ is
continuous. In particular, if $B=A$ and $\varphi$ is the identity, we see that
for $S\subset T$ we have a continuous homomorphism
$\rho^{T,S}:A\{S^{-1}\}\to A\{T^{-1}\}$ obtained by passing to the separated completion
from $S^{-1}A\to T^{-1}A$; if $U$ is a third multiplicative
subset of $A$ such that $S\subset T\subset U$, we have
$\rho^{U,S}=\rho^{U,T}\circ\rho^{T,S}$.
\end{env}

\begin{env}[7.6.8]
\label{0.7.6.8}
Let $S_1$, $S_2$ be two multiplicative subsets of $A$, and let $S_2'$ be the
canonical image of $S_2$ in $A\{S_1^{-1}\}$; we then have a canonical topological
isomorphism
$A\{(S_1S_2)^{-1}\}\xrightarrow{\sim}A\{S_1^{-1}\}\{{S_2'}^{-1}\}$,
as one sees by starting from the canonical isomorphism
$(S_1S_2)^{-1}A\xrightarrow{\sim}{S_2''}^{-1}(S_1^{-1}A)$ (where $S_2''$ is
the canonical image of $S_2$ in $S_1^{-1}A$), which is bicontinuous.
\end{env}

\begin{env}[7.6.9]
\label{0.7.6.9}
Let $\mathfrak a$ be an \emph{open} ideal of $A$; we may assume that
$\mathfrak J_\lambda\subset\mathfrak a$ for every $\lambda$, and consequently
$S^{-1}\mathfrak J_\lambda\subset S^{-1}\mathfrak a$ in the ring
$S^{-1}A$; in other words, $S^{-1}\mathfrak a$ is an \emph{open} ideal of
$S^{-1}A$; we shall denote by $\mathfrak a\{S^{-1}\}$ its separated completion,
equal to $\varprojlim(S^{-1}\mathfrak a/S^{-1}\mathfrak J_\lambda)$, which is an
\emph{open} ideal of $A\{S^{-1}\}$, isomorphic to the closure of the canonical
image of $S^{-1}\mathfrak a$. Moreover, \emph{the discrete ring
$A\{S^{-1}\}/\mathfrak a\{S^{-1}\}$ is canonically isomorphic to
$S^{-1}A/S^{-1}\mathfrak a=S^{-1}(A/\mathfrak a)$}. Conversely, if
$\mathfrak a'$ is an open ideal of $A\{S^{-1}\}$, $\mathfrak a'$ contains
an ideal of the form $\mathfrak J_\lambda\{S^{-1}\}$, hence is the inverse
image of an ideal of $S^{-1}A/S^{-1}\mathfrak J_\lambda$, which is
necessarily (1.2.6) of the form $S^{-1}\mathfrak a$, where
$\mathfrak a\supset\mathfrak J_\lambda$. We conclude that we have
$\mathfrak a'=\mathfrak a\{S^{-1}\}$. In particular (1.2.6):
\end{env}

\begin{proposition}[7.6.10]
\label{0.7.6.10}
The map $\mathfrak p\mapsto\mathfrak p\{S^{-1}\}$ is an increasing
bijection of the set of open prime ideals $\mathfrak p$ of $A$ such
that $\mathfrak p\cap S=\varnothing$ onto the set of open prime ideals

\oldpage[0\textsubscript{I}]{74}

of $A\{S^{-1}\}$; moreover, the field of fractions of
$A\{S^{-1}\}/\mathfrak p\{S^{-1}\}$ is canonically isomorphic to that of
$A/\mathfrak p$.
\end{proposition}

\begin{proposition}[7.6.11]
\label{0.7.6.11}
\begin{enumerate}
  \item[\rm(i)] If $A$ is an admissible ring, so is
    $A'=A\{S^{-1}\}$, and for every ideal of definition $\mathfrak J$ of $A$,
    $\mathfrak J'=\mathfrak J\{S^{-1}\}$ is an ideal of definition of $A'$.
  \item[\rm(ii)] Let $A$ be an adic ring and $\mathfrak J$ an ideal of
    definition of $A$ such that $\mathfrak J/\mathfrak J^2$ is of finite type
    over $A/\mathfrak J$; then $A'$ is a $\mathfrak J'$-adic ring and
    $\mathfrak J'/\mathfrak J'^2$ is of finite type over
    $A'/\mathfrak J'$. If moreover $A$ is Noetherian, so is
    $A'$.
\end{enumerate}
\end{proposition}

\begin{enumerate}
  \item[\rm(i)] If $\mathfrak J$ is an ideal of definition in $A$, it is
    clear that $S^{-1}\mathfrak J$ is an ideal of definition in the topological
    ring $S^{-1}A$, since we have
    $(S^{-1}\mathfrak J)^n=S^{-1}\mathfrak J^n$. Let $A''$ be the separated ring
    associated with $S^{-1}A$, and $\mathfrak J''$ the image of $S^{-1}\mathfrak J$ in
    $A''$; the image of $S^{-1}\mathfrak J^n$ is $\mathfrak J''^n$, so
    $\mathfrak J''^n$ tends to 0 in $A''$; since $\mathfrak J'$ is
    the closure of $\mathfrak J''$ in $A'$, $\mathfrak J'^n$ is contained
    in the closure of $\mathfrak J''^n$, and therefore tends to 0 in $A'$.
  \item[\rm(ii)] Set $A_i=A/\mathfrak J^{i+1}$, and for $i\leq j$, let
    $u_{ij}$ be the canonical homomorphism
    $A/\mathfrak J^{j+1}\to A/\mathfrak J^{i+1}$; let $S_i$ be the canonical
    image of $S$ in $A_i$, and set $A_i'=S_i^{-1}A_i$; finally, let
    $u_{ij}':A_j'\to A_i'$ be the homomorphism deduced canonically from $u_{ij}$.
    Let us show that the projective system $(A_i',u_{ij}')$ satisfies the conditions
    of Proposition~(7.2.7): it is clear that the $u_{ij}'$ are surjective;
    on the other hand, the kernel of $u_{ij}'$ is
    $S_j^{-1}(\mathfrak J^{i+1}/\mathfrak J^{j+1})$ (1.3.2), equal to
    $\mathfrak J_j'^{i+1}$, where
    $\mathfrak J_j'=S_j^{-1}(\mathfrak J/\mathfrak J^{j+1})$; finally,
    $\mathfrak J_1'/\mathfrak J_1'^2=S_1^{-1}(\mathfrak J/\mathfrak J^2)$,
    and since $\mathfrak J/\mathfrak J^2$ is of finite type over
    $A/\mathfrak J^2$, $\mathfrak J_1'/\mathfrak J_1'^2$ is of finite type
    over $A_1'$. Finally, if $A$ is Noetherian, so is
    $A_0'=S_0^{-1}(A/\mathfrak J)$, which, by Corollary (7.2.8), completes the proof of Proposition
    (7.6.11).
\end{enumerate}

\begin{corollary}[7.6.12]
\label{0.7.6.12}
Under the hypotheses of (7.6.11, (ii)), we have
$(\mathfrak J\{S^{-1}\})^n=\mathfrak J^n\{S^{-1}\}$.
\end{corollary}

This follows indeed from (7.2.7) and from the proof of (7.6.11).

\begin{proposition}[7.6.13]
\label{0.7.6.13}
Let $A$ be a Noetherian adic ring and $S$ a multiplicative subset of $A$;
then $A\{S^{-1}\}$ is a flat $A$-module.
\end{proposition}

Indeed, if $\mathfrak J$ is an ideal of definition of $A$,
$A\{S^{-1}\}$ is the separated completion of the Noetherian ring $S^{-1}A$ endowed
with the $S^{-1}\mathfrak J$-preadic topology; consequently (7.3.3)
$A\{S^{-1}\}$ is a flat $S^{-1}A$-module; since $S^{-1}A$ is a
flat $A$-module (6.3.1), the proposition follows from the transitivity of
flatness (6.2.1).

\begin{corollary}[7.6.14]
\label{0.7.6.14}
Under the hypotheses of (7.6.13), let $S'\subset S$ be a second multiplicative
subset of $A$; then $A\{S^{-1}\}$ is a flat
$A\{{S'}^{-1}\}$-module.
\end{corollary}

Indeed (7.6.8), $A\{S^{-1}\}$ is canonically identified with
$A\{{S'}^{-1}\}\{S_0^{-1}\}$, where $S_0$ is the canonical image of $S$ in
$A\{{S'}^{-1}\}$, and $A\{{S'}^{-1}\}$ is Noetherian (7.6.11).

\begin{env}[7.6.15]
\label{0.7.6.15}
For every element $f$ of a linearly topologized ring $A$, we shall
denote by $A_{\{f\}}$ the complete ring of fractions $A\{S_f^{-1}\}$,
where $S_f$ is the multiplicative set of the $f^n$ ($n\geq 0$); for every
open ideal $\mathfrak a$ of $A$, we shall write $\mathfrak a_{\{f\}}$ instead
of $\mathfrak a\{S_f^{-1}\}$. If $g$ is a second element of $A$, we have
a continuous canonical homomorphism $A_{\{f\}}\to A_{\{fg\}}$ (7.6.7).
When $f$ runs through a multiplicative subset $S$ of $A$, the $A_{\{f\}}$
therefore form a filtered inductive system of rings for the preceding
homomorphisms; we shall set
$A_{\{S\}}=\varinjlim_{f\in S}A_{\{f\}}$. For every $f\in S$, we have a
homomorphism $A_{\{f\}}\to A\{S^{-1}\}$ (7.6.7), and

\oldpage[0\textsubscript{I}]{75}

these homomorphisms form an inductive system; by passage to the inductive
limit, they therefore define a canonical homomorphism
$A_{\{S\}}\to A\{S^{-1}\}$.
\end{env}

\begin{proposition}[7.6.16]
\label{0.7.6.16}
If $A$ is a Noetherian ring, $A\{S^{-1}\}$ is a flat module over
$A_{\{S\}}$.
\end{proposition}

Indeed (7.6.14), $A\{S^{-1}\}$ is flat over each of the rings
$A_{\{f\}}$ for $f\in S$, and the conclusion follows from (6.2.3).

\begin{proposition}[7.6.17]
\label{0.7.6.17}
Let $\mathfrak p$ be an open prime ideal in an admissible ring $A$, and
let $S=A-\mathfrak p$. Then the rings $A\{S^{-1}\}$ and $A_{\{S\}}$ are
local rings, the canonical homomorphism
$A_{\{S\}}\to A\{S^{-1}\}$ is local, and the residue fields of $A_{\{S\}}$
and $A\{S^{-1}\}$ are canonically isomorphic to the field of fractions of
$A/\mathfrak p$.
\end{proposition}

Indeed, let $\mathfrak J\subset\mathfrak p$ be an ideal of definition of
$A$; we have $S^{-1}\mathfrak J\subset S^{-1}\mathfrak p=\mathfrak p
A_{\mathfrak p}$, so $A_{\mathfrak p}/S^{-1}\mathfrak J$ is a local
ring; we conclude from (7.1.12), (7.6.9) and (7.6.11, (i)) that
$A\{S^{-1}\}$ is a local ring. Set
$\mathfrak m=\varinjlim_{f\in S}\mathfrak p_{\{f\}}$, which is an ideal of
$A_{\{S\}}$; we shall see that every element of $A_{\{S\}}$ not belonging
to $\mathfrak m$ is invertible. Indeed, such an element is the image in
$A_{\{S\}}$ of an element $z\in A_{\{f\}}$ not belonging to
$\mathfrak p_{\{f\}}$, for some $f\in S$; its canonical image $z_0$ in
$A_{\{f\}}/\mathfrak J_{\{f\}}=S_f^{-1}(A/\mathfrak J)$ therefore does not belong
to $S_f^{-1}(\mathfrak p/\mathfrak J)$ (7.6.9), which means that
$z_0=\overline{x}/\overline{f}^k$, where $x\notin\mathfrak p$ and
$\overline{x},\overline{f}$ are the classes of $x,f$ modulo~$\mathfrak J$.
Since $x\in S$, we have $g=xf\in S$, and in $S_g^{-1}A$,
$y_0=x^{k+1}/g^k$, the canonical image of $x/f^k\in S_f^{-1}A$, has an
inverse $x^{k-1}f^{2k}/g^k$. This implies \emph{a fortiori} that the image of
$y_0$ in $S_g^{-1}A/S_g^{-1}\mathfrak J$ is invertible, hence (7.6.9 and
7.1.12) the canonical image $y$ of $z$ in $A_{\{g\}}$ is invertible;
the image of $z$ in $A_{\{S\}}$ (equal to that of $y$) is consequently
invertible. We therefore see that $A_{\{S\}}$ is a local ring with maximal ideal
$\mathfrak m$; moreover, the image of $\mathfrak p_{\{f\}}$ in
$A\{S^{-1}\}$ is contained in the maximal ideal $\mathfrak p\{S^{-1}\}$ of
this ring; \emph{a fortiori}, the image of $\mathfrak m$ in
$A\{S^{-1}\}$ is contained in $\mathfrak p\{S^{-1}\}$, so
the canonical homomorphism $A_{\{S\}}\to A\{S^{-1}\}$ is local. Finally,
since every element of $A\{S^{-1}\}/\mathfrak p\{S^{-1}\}$ is the image of an
element of a ring $S_f^{-1}A$ for a suitable $f\in S$,
the homomorphism $A_{\{S\}}\to A\{S^{-1}\}/\mathfrak p\{S^{-1}\}$ is
surjective, and therefore gives, by passage to quotients, an isomorphism of the residue
fields.

\begin{corollary}[7.6.18]
\label{0.7.6.18}
Under the hypotheses of (7.6.17), if we suppose moreover that $A$ is a Noetherian
adic ring, the local rings $A\{S^{-1}\}$ and $A_{\{S\}}$ are
Noetherian, and $A\{S^{-1}\}$ is a faithfully flat $A_{\{S\}}$-module.
\end{corollary}

We already know (7.6.11, (ii)) that $A\{S^{-1}\}$ is Noetherian and
$A_{\{S\}}$-flat (7.6.16); since the homomorphism
$A_{\{S\}}\to A\{S^{-1}\}$ is local, we conclude that $A\{S^{-1}\}$ is
a faithfully flat $A_{\{S\}}$-module (6.6.2), and consequently that $A_{\{S\}}$
is Noetherian (6.5.2).

\subsection{Completed tensor products}
\label{subsection:0.7.7}

\begin{env}[7.7.1]
\label{0.7.7.1}
Let $A$ be a linearly topologized ring and $M$, $N$ two linearly
topologized $A$-modules. Let $\mathfrak J$, $V$, $W$ be open neighbourhoods
of $0$ in $A$, $M$, $N$ respectively, which are
$A$-modules and are such that $\mathfrak JM\subset V$, $\mathfrak JN\subset W$,
so that $M/V$ and $N/W$ can be considered as
$(A/\mathfrak J)$-modules. When $\mathfrak J$, $V$, $W$ run through the
systems of open neighbourhoods satisfying the preceding conditions, it is
immediate that the modules
$(M/V)\otimes_{A/\mathfrak J}(N/W)$ form a projective system of modules

\oldpage[0\textsubscript{I}]{76}

over the projective system of rings $A/\mathfrak J$; by passage to the projective
limit, one therefore deduces from them a module over the separated completion $\widehat A$
of $A$, which is called the \emph{completed tensor product} of $M$ and
$N$ and is denoted $(M\otimes_A N)^\wedge$. If one notes that $M/V$
is canonically isomorphic to $\widehat M/\overline V$, where $\widehat M$ is
the separated completion of $M$ and $\overline V$ the closure in $\widehat M$ of
the image of $V$, one sees that the completed tensor product
$(M\otimes_A N)^\wedge$ is canonically identified with
$(\widehat M\otimes_{\widehat A}\widehat N)^\wedge$, which is also denoted
$\widehat M\widehat\otimes_{\widehat A}\widehat N$.
\end{env}

\begin{env}[7.7.2]
\label{0.7.7.2}
\begin{sloppypar}
With the preceding notation, the tensor products
$(M/V)\otimes_A(N/W)$ and
$(M/V)\otimes_{A/\mathfrak J}(N/W)$ are canonically identified; they
are therefore also identified with
$(M\otimes_A N)/(\Im(V\otimes_A N)+\Im(M\otimes_A W))$.
We conclude that $(M\otimes_A N)^\wedge$ is \emph{the separated completion of the
$A$-module $M\otimes_A N$, endowed with the topology for which the
submodules}
\[
  \Im(V\otimes_A N)+\Im(M\otimes_A W)
\]
\emph{form a fundamental system of neighbourhoods of $0$} ($V$ and $W$
running through the set of open submodules of $M$ and $N$
respectively); we shall say for short that this topology is the
\emph{tensor product} of the given topologies on $M$ and $N$.
\end{sloppypar}
\end{env}

\begin{env}[7.7.3]
\label{0.7.7.3}
Let $M'$, $N'$ be two linearly topologized $A$-modules and
$u:M\to M'$, $v:N\to N'$ two continuous homomorphisms; it is immediate
that $u\otimes v$ is continuous for the tensor product topologies on
$M\otimes_A N$ and $M'\otimes_A N'$ respectively; by passage to the separated
completion, one deduces from it a continuous homomorphism
$(M\otimes_A N)^\wedge\to(M'\otimes_A N')^\wedge$, which we shall denote
by $u\widehat\otimes v$; $(M\otimes_A N)^\wedge$ is therefore a
\emph{bifunctor} in $M$ and $N$ on the category of linearly topologized
$A$-modules.
\end{env}

\begin{env}[7.7.4]
\label{0.7.7.4}
One defines in the same way the completed tensor product of an arbitrary finite number
of linearly topologized $A$-modules; it is immediate that this product
has the usual properties of associativity and commutativity.
\end{env}

\begin{env}[7.7.5]
\label{0.7.7.5}
If $B$, $C$ are two linearly topologized topological $A$-algebras,
the tensor product topology on $B\otimes_A C$ has as a fundamental system
of neighbourhoods of $0$ the \emph{ideals}
$\Im(\mathfrak K\otimes_A C)+\Im(B\otimes_A\mathfrak L)$ of the algebra
$B\otimes_A C$, $\mathfrak K$ (resp. $\mathfrak L$) running through the set
of open ideals of $B$ (resp. $C$). Consequently, $(B\otimes_A C)^\wedge$
is endowed with a structure of \emph{topological $\widehat A$-algebra}, the projective
limit of the projective system of $(A/\mathfrak J)$-algebras
$(B/\mathfrak K)\otimes_{A/\mathfrak J}(C/\mathfrak L)$ ($\mathfrak J$
an open ideal of $A$ such that $\mathfrak J.B\subset\mathfrak K$,
$\mathfrak J.C\subset\mathfrak L$; such ideals always exist). This
algebra is called the \emph{completed tensor product} of the algebras $B$ and
$C$.
\end{env}

\begin{env}[7.7.6]
\label{0.7.7.6}
The homomorphisms of $A$-algebras $b\mapsto b\otimes 1$,
$c\mapsto 1\otimes c$ from $B$ and $C$ to $B\otimes_A C$ are continuous
when the latter algebra is endowed with the tensor product topology;
by composition with the canonical homomorphism from $B\otimes_A C$ to its
separated completion, they therefore give \emph{canonical homomorphisms}
$\rho:B\to(B\otimes_A C)^\wedge$, $\sigma:C\to(B\otimes_A C)^\wedge$.
The algebra $(B\otimes_A C)^\wedge$ and the homomorphisms $\rho$ and $\sigma$
moreover have the \emph{following universal property}: for every
linearly topologized, separated and complete $A$-algebra $D$, and every
pair of continuous $A$-homomorphisms $u:B\to D$, $v:C\to D$, there exists one
and only one continuous $A$-homomorphism
$w:(B\otimes_A C)^\wedge\to D$ such that $u=w\circ\rho$ and
$v=w\circ\sigma$. Indeed, there already exists a unique $A$-homomorphism
$w_0:B\otimes_A C\to D$ such that $u(b)=w_0(b\otimes 1)$ and
$v(c)=w_0(1\otimes c)$, and everything comes down to proving that $w_0$ is continuous,
because it will then give a continuous

\oldpage[0\textsubscript{I}]{77}

homomorphism $w$ by passage to the separated completion. Now, if
$\mathfrak M$ is an open ideal of $D$, there exist by hypothesis open
ideals $\mathfrak K\subset B$, $\mathfrak L\subset C$ such that
$u(\mathfrak K)\subset\mathfrak M$, $v(\mathfrak L)\subset\mathfrak M$;
the image under $w_0$ of
$\Im(\mathfrak K\otimes_A C)+\Im(B\otimes_A\mathfrak L)$ is again
contained in $\mathfrak M$, whence our assertion.
\end{env}

\begin{proposition}[7.7.7]
\label{0.7.7.7}
If $B$ and $C$ are two preadmissible $A$-algebras,
$(B\otimes_A C)^\wedge$ is admissible, and if $\mathfrak K$ (resp.
$\mathfrak L$) is an ideal of definition of $B$ (resp. $C$), the closure
in $(B\otimes_A C)^\wedge$ of the canonical image of
$\mathfrak H=\Im(\mathfrak K\otimes_A C)+\Im(B\otimes_A\mathfrak L)$ is
an ideal of definition.
\end{proposition}

It suffices to show that $\mathfrak H^n$ tends to $0$ for the tensor
product topology, which follows immediately from the inclusion
\[
  \mathfrak H^{2n}\subset
  \Im(\mathfrak K^n\otimes_A C)+\Im(B\otimes_A\mathfrak L^n).
\]

\begin{proposition}[7.7.8]
\label{0.7.7.8}
Let $A$ be a preadic ring, $\mathfrak J$ an ideal of definition of
$A$, and $M$ an $A$-module \emph{of finite type}, endowed with the
$\mathfrak J$-preadic topology. For every adic and Noetherian topological
$A$-algebra $B$, $B\otimes_A M$ is identified with the completed tensor product
$(B\otimes_A M)^\wedge$.
\end{proposition}

If $\mathfrak K$ is an ideal of definition of $B$, there exists by hypothesis
an integer $m$ such that $\mathfrak J^mB\subset\mathfrak K$, hence
$\Im(B\otimes_A\mathfrak J^{nm}M)=
\Im(\mathfrak J^{nm}B\otimes_A M)\subset
\Im(\mathfrak K^nB\otimes_A M)=\mathfrak K^n(B\otimes_A M)$; we
conclude that on $B\otimes_A M$, the tensor product of the topologies of $B$
and of $M$ is the $\mathfrak K$-preadic topology. Since $B\otimes_A M$
is a $B$-module of finite type, the proposition follows at once from (7.3.6).

\subsection{Topologies on modules of homomorphisms}
\label{subsection:0.7.8}

\begin{env}[7.8.1]
\label{0.7.8.1}
Let $A$ be a Noetherian $\mathfrak J$-adic ring, and $M$ and $N$ two
$A$-modules of finite type, endowed with the $\mathfrak J$-preadic topology;
we know (7.3.6) that they are separated and complete; moreover, every
$A$-homomorphism $M\to N$ is automatically continuous, and the $A$-module
$\operatorname{Hom}_A(M,N)$ is of finite type. For every integer $i\geq 0$, set
$A_i=A/\mathfrak J^{i+1}$, $M_i=M/\mathfrak J^{i+1}M$,
$N_i=N/\mathfrak J^{i+1}N$; for $i\leq j$, every homomorphism
$u_j:M_j\to N_j$ maps $\mathfrak J^{i+1}M_j$ into
$\mathfrak J^{i+1}N_j$, and hence gives, by passage to quotients, a
homomorphism $u_i:M_i\to N_i$, which defines a canonical homomorphism
$\operatorname{Hom}_{A_j}(M_j,N_j)\to\operatorname{Hom}_{A_i}(M_i,N_i)$; moreover, the
$\operatorname{Hom}_{A_i}(M_i,N_i)$ constitute a \emph{projective system} for these
homomorphisms, and it follows from (7.2.10) that there is a canonical
isomorphism
\[
  \varphi:\operatorname{Hom}_A(M,N)\longrightarrow
  \varprojlim_i\operatorname{Hom}_{A_i}(M_i,N_i).
\]
Moreover:
\end{env}

\begin{proposition}[7.8.2]
\label{0.7.8.2}
If $M$ and $N$ are modules of finite type over a Noetherian $\mathfrak J$-adic
ring $A$, the submodules $\operatorname{Hom}_A(M,\mathfrak J^{i+1}N)$ form a
fundamental system of neighbourhoods of $0$ in $\operatorname{Hom}_A(M,N)$ for the
$\mathfrak J$-adic topology, and the canonical homomorphism
\[
  \varphi:\operatorname{Hom}_A(M,N)\longrightarrow
  \varprojlim_i\operatorname{Hom}_{A_i}(M_i,N_i)
\]
is a topological isomorphism.
\end{proposition}

Indeed, we can consider $M$ as a quotient of a free $A$-module $L$ of
finite type, and consequently identify $\operatorname{Hom}_A(M,N)$ with a submodule of
$\operatorname{Hom}_A(L,N)$; under this identification,
$\operatorname{Hom}_A(M,\mathfrak J^{i+1}N)$ is the intersection of $\operatorname{Hom}_A(M,N)$ and
$\operatorname{Hom}_A(L,\mathfrak J^{i+1}N)$; since the topology induced on
$\operatorname{Hom}_A(M,N)$ by the $\mathfrak J$-adic topology of $\operatorname{Hom}_A(L,N)$ is
the $\mathfrak J$-adic topology (7.3.2), we are reduced to proving the
first assertion

\oldpage[0\textsubscript{I}]{78}

for $M=L=A^m$; but then $\operatorname{Hom}_A(L,N)=N^m$,
$\operatorname{Hom}_A(L,\mathfrak J^{i+1}N)=(\mathfrak J^{i+1}N)^m=
\mathfrak J^{i+1}N^m$, and the result is obvious. To establish the second
assertion, note that the image of $\operatorname{Hom}_A(M,\mathfrak J^{i+1}N)$ in
$\operatorname{Hom}_{A_j}(M_j,N_j)$ is zero for $j\leq i$, so $\varphi$ is
continuous; conversely, the inverse image in $\operatorname{Hom}_A(M,N)$ of the $0$ of
$\operatorname{Hom}_{A_i}(M_i,N_i)$ is $\operatorname{Hom}_A(M,\mathfrak J^{i+1}N)$, so $\varphi$
is bicontinuous.

If we suppose only that $A$ is a Noetherian $\mathfrak J$-preadic
ring, and $M$ and $N$ two $A$-modules of finite type that are separated for the
$\mathfrak J$-preadic topology, the preceding proof shows that
the first assertion of (7.8.2) remains valid, and that $\varphi$ is a
topological isomorphism of $\operatorname{Hom}_A(M,N)$ onto a submodule of
$\varprojlim_i\operatorname{Hom}_{A_i}(M_i,N_i)$.

\begin{proposition}[7.8.3]
\label{0.7.8.3}
Under the hypotheses of (7.8.2), the set of injective (resp. surjective,
bijective) homomorphisms from $M$ to $N$ is an open subset of
$\operatorname{Hom}_A(M,N)$.
\end{proposition}

Indeed, by virtue of (7.3.5) and (7.1.14), for $u$ to be surjective, it is
necessary and sufficient that the corresponding homomorphism
$u_0:M/\mathfrak JM\to N/\mathfrak JN$ be so, and the set of
surjective homomorphisms from $M$ to $N$ is therefore the inverse image under
the continuous map $\operatorname{Hom}_A(M,N)\to\operatorname{Hom}_{A_0}(M_0,N_0)$ of a subset
of a discrete space. Let us now show that the set of injective homomorphisms
is open; let $v$ be such a homomorphism, and set $M'=v(M)$;
by the Artin--Rees lemma (7.3.2.1), there exists an integer $k\geq 0$ such that
$M'\cap\mathfrak J^{m+k}N\subset\mathfrak J^mM'$ for every $m>0$; we
shall see that for $w\in\mathfrak J^{k+1}\operatorname{Hom}_A(M,N)$,
$u=v+w$ is injective, which will complete the proof. Indeed, let
$x\in M$ be such that $u(x)=0$; let us prove that for every $i\geq 0$ the relation
$x\in\mathfrak J^iM$ implies $x\in\mathfrak J^{i+1}M$; it will indeed
follow that
\[
  x\in\bigcap_{i\geq 0}\mathfrak J^iM=(0).
\]
Indeed, we then have $w(x)\in\mathfrak J^{i+k+1}N$, and moreover
$w(x)=-v(x)\in M'$, hence
$v(x)\in M'\cap\mathfrak J^{i+k+1}N\subset\mathfrak J^{i+1}M'$, and since
$v$ is an isomorphism of $M$ onto $M'$, $x\in\mathfrak J^{i+1}M$;
Q.E.D.

\begin{flushright}
\emph{(To be continued.)}
\end{flushright}

